%% kapa: exam=1 round=24 year=2013 period=1 form=A subject=civil_law subject_ko=민법 seq=18 number=18
%% source: Q-Net 감정평가사 기출문제 (cst00309) artlSeq=5073802 file=extracted/1008/5073802/2013년 제24회 감정평가사 1차시험 기출문제/1교시(A형).hwp
%% conversion: kordoc 4.24.51 -> tools/gen1.py (auto); review=proofread-sonnet; flags=none
\begin{question}{18}
소멸시효에 관한 설명으로 옳지 않은 것은? (다툼이 있으면 판례에 의함)
\begin{choices}
  \item 소멸시효의 기간은 법률행위로 단축할 수 있다.
  \item 매달 지급해야 하는 집합건물의 관리비채권의 소멸시효기간은 10년이다.
  \item 환매권의 행사로 발생한 소유권이전등기청구권의 소멸시효기간은 10년이다.
  \item 공유물분할청구권은 그 기초가 되는 법률관계가 존속하는 동안 시효로 소멸하지 않는다.
  \item 근저당설정약정에 의한 근저당권설정등기청구권은 피담보채권과 별도로 소멸시효가 진행한다.
\end{choices}
\end{question}
\answer{2}
